% Compiles every style of the dzg TikZ library through the real dzg-tikz
% package, so a renamed, removed, or broken style fails `just test`.
\documentclass{article}
\usepackage{dzg-tikz}
\pagestyle{empty}

\begin{document}

\section*{Coxeter diagrams}
\begin{tikzpicture}
  \foreach \m [count=\k] in {3,4,5,6,infinity,dotted} {
    \node[vertex=white] (a\k) at (0,-\k) {};
    \node[vertex=black] (b\k) at (1.5,-\k) {};
    \begin{scope}[on edge layer]
      \draw[coxeter edge=\m] (a\k.center) -- (b\k.center);
    \end{scope}
  }
  \node[vertex=doubled, label={[vertex label]above:$r_0$}] (c) at (3,-1) {};
  \node[vertex=white] (d) at (4.5,-1) {};
  \draw[coxeter edge label=8] (c) -- (d);
  \scoped[on background layer] \draw[parabolic] (c.center) -- (d.center);
  \draw[fold] (3,-2) -- (4.5,-2);
  \draw[fold axis] (3,-3) -- (4.5,-3);
  \node[reflected root] at (b1) {};
  \draw[quotient map] (3,-4) -- (4.5,-4);
\end{tikzpicture}

\section*{Coxeter constructors}
The keys of the Coxeter constructors (label schemes, parity marks, maximal
parabolics, mirror targets) on objects, and the generic pics.

\begin{tikzpicture}[scale=0.4]
  \foreach \s [count=\k] in {alpha, r, ell, index, none}
    \pic[root labels=\s] at (6*\k,0) {object=coxeter/vinberg-18-2-0};
  \pic[root labels, parity marks=false] at (36,0) {object=coxeter/vinberg-18-2-0};
  \pic[root labels=index] (K) at (0,-10) {object=coxeter/sterk-cusp-cover-3};
  \scoped[on background layer] \draw[parabolic] (K-0.center) -- (K-20.center);
  \pic (A) [maximal parabolic=B8] at (8,-10) {object=coxeter/vinberg-10-8-0};
  \pic (V) [maximal parabolic=D16] at (20,-10) {object=coxeter/vinberg-18-0-0};
  \pic (W) [maximal parabolic=E7+D10] at (30,-10) {object=coxeter/vinberg-19-1-1};
  \pic[mirror targets=all] at (0,-22) {object=mirror-moves/coble};
  \pic at (12,-22) {dynkin folding={ply=2}{E{oooooo}}{}{F{**oo}}};
  \pic (R) at (24,-22) {root chain={black/4/3, black/4/4, white/2/}};
  \draw[fold] (R-0) -- (R-2);
\end{tikzpicture}

\coxeterpair\quad \coxeterpair[marks=black/doubled, edge=none, roots=f_1/f_2]\quad
\coxeterpair[edge=dotted, weight=$h_1h_2$]\quad
\coxetervertex{white}\coxetervertex{black}\coxetervertex{doubled}

\section*{Baily--Borel cusp diagrams}
\begin{tikzpicture}
  \node[cusp0, label={[cusp label]above:$\delta=0$}] (p) at (0,0) {$U\oplus E_8$};
  \node[cusp1] (q) at (0,-1.5) {$E_8^2$};
  \draw[incidence] (p) -- (q);
  \draw[cusp brace] (-1,-1.5) -- (1,-1.5);
\end{tikzpicture}

The keys of the cusp-diagram constructors: each label scheme and incidence
divisibility, on objects that use \verb|\bblabel|, \verb|\bbbrace| and
\verb|\bbcolumn|.

\begin{tikzpicture}[scale=0.45]
  \pic (S) {object=bb-cusps/fen2};
  \pic[cusp labels=eta, incidence divisibility] at (22,0) {object=bb-cusps/fen2};
  \pic[cusp labels=eta] at (0,-18) {object=bb-cusps/f220};
  \pic[cusp labels=none] at (22,-18) {object=bb-cusps/fen};
  \pic[cusp labels=none] at (0,-30) {object=bb-cusps/f2d-5};
  \scoped[on background layer] \draw[cusp image] (S-2.center) -- (S-245.center);
\end{tikzpicture}

\section*{IAS, fans, Kulikov models}
\begin{tikzpicture}
  \fill[ias region] (0,0) rectangle (2,2);
  \draw[ias boundary] (0,0) rectangle (2,2);
  \draw[ias edge] (0,0) -- (2,2);
  \node[ias singularity=16] at (1,1.5) {};
  \node[ias singularity] at (1.5,0.5) {};
  \filldraw[ias surgery] (0.2,0.2) -- (0.6,0.2) -- (0.4,0.5) -- cycle;
  \draw[fan ray] (3,1) -- (4,1.5);
  \node[kulikov component, minimum width=1cm, minimum height=1.5cm] at (5,1) {$V_1$};
  \draw[double curve] (5.5,0.25) -- (5.5,1.75);
  \node[nontoric blowup] at (6.5,1) {};
  \node[triple point] at (7.5,1) {};
  \node[self intersection] at (8,1) {$-2$};
\end{tikzpicture}

The keys of the polygon constructors; the eta\_3 instance collapses vertices
and merges two surgery points.

\begin{tikzpicture}[scale=0.3]
  \pic[ias labels=multiplicities, ell={2,0,0,0,0,0,0,0,0,0,0,0,0,0,0,0,2,4,6,4}]
    (A) {object=ias/symington-18-2-0};
  \pic[ias labels=none] at (16,0) {object=ias/symington-18-0-0};
  \pic[ias involutions] at (5,-20) {object=ias/sterk3-ias};
  \draw[fold] (A-p16) -- (A-p18);
\end{tikzpicture}

\begin{tikzpicture}[scale=0.5]
  \pic[ade labels=none, ade title=false] (Q) {object=ade/minus-E6-minus};
  \draw[fold] (Q-pstar) -- (Q-1);
  \pic[fan labels=none, fan length=1] (F) at (10,0) {object=toric/fan-F2};
  \draw[polytope edge, dashed] (F-rho1) -- (F-rho2) -- (F-rho3) -- (F-rho4) -- cycle;
\end{tikzpicture}

\section*{Objects}
Every file under figures/objects, placed through the object pic; the list is
passed in \verb|\dzgobjects| by the test recipe.


\foreach \o in \dzgobjects {\par\noindent\texttt{\o}\par
  \begin{tikzpicture}[scale=0.4]\pic {object=\o};\end{tikzpicture}\par}
\section*{Lattice polygons}
\begin{tikzpicture}
  \draw[lattice grid] (-0.5,-0.5) grid (3.5,2.5);
  \fill[polygon region] (0,0) -- (3,0) -- (0,2) -- cycle;
  \draw[long side] (0,0) -- (3,0);
  \draw[short side] (3,0) -- (0,2) -- (0,0);
  \node[lattice point] at (1,1) {};
  \node[boundary point] at (1,0) {};
  \node[marked point=p] at (2,0) {};
\end{tikzpicture}

\section*{Posets and moduli schematics}
\begin{tikzpicture}
  \node[poset node] (x) at (0,0) {$\overline{M}_2$};
  \node[poset node] (y) at (0,-1.5) {$\Delta_0$};
  \draw[degeneration] (x) -- (y);
  \draw[moduli blob] (2,0) to[closed, curve through={(3,0.4) (3.5,-0.5) (2.5,-1)}] (2,0);
  \node[stratum] at (2.7,-0.3) {};
  \draw[boundary stratum] (2.2,-0.8) -- (3.2,-0.8);
\end{tikzpicture}

\section*{Nikulin's table of 2-elementary lattices}
\begin{tikzpicture}[x=0.3cm, y=0.3cm]
  \pic (N) {object=nikulin/two-elementary-table};
  \pic[lattice table labels=none] at (24,0) {object=nikulin/two-elementary-table};
  \scoped[on background layer] \draw[parabolic] (N-2-2.center) -- (N-10-10.center);
\end{tikzpicture}

\section*{Strata of $\overline{M}_2$}
\begin{tikzpicture}[x=0.5cm, y=0.5cm]
  \pic[genus labels=none] (C) {object=moduli/m2-curve-4};
  \pic (X) at (26,0) {object=moduli/m2-coxeter};
  \pic[genus labels=none] (V) at (26,-6) {object=moduli/m2-graph-6};
  \draw[degeneration] (X-c) -- (V-box);
\end{tikzpicture}

\section*{Stable curves and dual graphs}
\begin{tikzpicture}
  \draw[curve component] (0,0) to[curve through={(0.6,0.4) (1.2,-0.2)}] (1.8,0.2);
  \node[curve node] at (0.6,0.4) {};
  \node[dual vertex, label={[genus]above:$1$}] at (3,0) {};
\end{tikzpicture}

\section*{Spaces, contours and spectral sequences}
Every constructor and key in constructors/dzg-spaces.tex: the space objects
without labels and decorated on their named points, the axes of C, and ss pages
in both conventions.

\begin{tikzpicture}[scale=0.5]
  \foreach \n in {1,2,3} \pic[space labels=none] at (3*\n,0) {object=spaces/ball-\n};
  \pic[space labels=none] (T) at (16,0) {object=spaces/torus};
  \draw[distinguished curve] (T-longitude) ellipse (3 and 1);
  \pic[scale=0.5, space labels=none] (S) at (22,-2) {object=spaces/torus-fundamental-square};
  \draw[lattice grid] (S-00) -- (S-11);
  \pic[space labels=none] (R) at (0,-8) {object=contours/strip-rectangle};
  \draw[branch cut] (R-mRb) -- (R-Rb);
  \begin{scope}[shift={(9,-8)}]\complexaxes{-1}{2}{-1}{2}\end{scope}
  \pic (E) at (16,-16) [ss homological, ss p={0,...,3}, ss q={0,1},
    ss q labels={0,n-1}] {ss page};
  \ssentry{E}{2}{0}{$H_2(B)$}
  \ssdifferential{E}{2}{2}{0}
  \pic (F) at (0,-24) [ss p={-1,...,1}, ss q={-1,...,1}, ss ticks=false] {ss page};
  \ssdifferential[ss unlabelled]{F}{2}{-1}{1}
\end{tikzpicture}

\end{document}
